\section{Matching width laws throughout the subcritical interval}\label{sec:rates54}
We now distinguish the dimension of the command arithmetic from the dimension of the physical orbit. Let $s\ge1$ and let the alphabet contain the rotations $R_{2\pi\alpha_1},\ldots,R_{2\pi\alpha_s}$ and $I$. For the matching upper bound the alphabet is exactly these letters, optionally with the reflection $F=\diag(1,-1)$. Use signed coordinate seeds of radius $\rho$ and coordinate queries. Extra letters do not affect the lower bounds, but are not silently included in the upper construction.

For a real number $x$, write $\|x\|_{\R/\Z}$ for distance to the nearest integer. The dual Diophantine condition is
\begin{equation}\label{eq:DA54}
 \|n\cdot\alpha\|_{\R/\Z}\ge b\|n\|_1^{-\tau}
 \qquad(0\ne n\in\Z^s),\qquad b>0,\quad\tau\ge s.
\end{equation}
In particular $1,\alpha_1,\ldots,\alpha_s$ are rationally independent. When $\tau=s$ this is the dual badly approximable condition. All implicit constants below may depend on $\alpha,b,s,\rho,\varepsilon$, but not on $N$ or a stochastic realization.

\begin{theorem}[Full-subcritical matching law]\label{thm:matched54}
Suppose \eqref{eq:DA54} holds with $\tau=s$, and let $0<\rho<1$. For every fixed $0\le\varepsilon<\rho/2$,
\begin{equation}\label{eq:matched54}
                 W_{N,\varepsilon}=\Theta\bigl(N^{s/(2s+1)}\bigr).
\end{equation}
The upper bound is achieved by an exact horizon-dependent realization. At $\varepsilon\ge\rho/2$, one command label suffices. If $\tau\ge s$ is arbitrary, the lower bound $W_{N,\varepsilon}=\Omega(N^{s/(2\tau+1)})$ still holds throughout the same subcritical interval.
\end{theorem}
The exponent for one badly approximable rotation is $1/3$, not logarithmic. The statement includes zero error when the signal has strict amplitude slack $\rho<1$. The threshold theorem continues to include $\rho=1$, but the exact polynomial upper bound here is not asserted at that endpoint. The earlier arithmetic argument obtained this scale in a smaller noise range. The proof below supplies the missing full-noise lower estimate and states the exact upper realization in full.

\begin{proposition}[An explicit family in every command dimension]\label{prop:algebraicangles54}
Let $\xi$ be a real algebraic integer of degree $s+1$. Then
$\alpha=(\xi,\xi^2,\ldots,\xi^s)$ satisfies \eqref{eq:DA54} with $\tau=s$ and some effectively computable $b>0$.
\end{proposition}
\begin{proof}
For $0\ne n\in\Z^s$ put $H=\|n\|_1$, and choose $n_0\in\Z$ so that
$z=n_0+\sum_{l=1}^s n_l\xi^l$ has modulus $\|n\cdot\alpha\|_{\R/\Z}$.
The degree assumption gives $z\ne0$. Write $\xi_1=\xi,\ldots,\xi_{s+1}$ for the conjugates and put $M_j=\max_{1\le l\le s}|\xi_j|^l$. Then
$|n_0|\le H M_1+1/2$ and
\[
 \left|n_0+\sum_{l=1}^s n_l\xi_j^l\right|
 \le H(M_1+M_j+1/2)\qquad(j\ge2).
\]
The norm of the nonzero algebraic integer $z$ is a nonzero integer. Thus
\[
 |z|\ge H^{-s}\prod_{j=2}^{s+1}(M_1+M_j+1/2)^{-1}.
\]
Choose a positive rational $b$ below this computable positive constant and $1/2$. This proves the claim. It concerns algebraic \emph{angles}, not algebraic entries of their rotation matrices.
\end{proof}

\subsection{Harmonic localization}
For the harmonic calculation assume that the alphabet contains $I,R_\theta$ and that one seed-query pair has mean $\rho\cos\phi$ after accumulated angle $\phi$, where $0<\rho\le1$. Other seeds, queries, or commands impose additional constraints and do not invalidate the lower bound. Let $\lambda=e^{i\theta}$ have infinite order. Put $\delta=2\varepsilon<\rho$, and choose an integer
\begin{equation}\label{eq:r53}
 r\ge1,\qquad r>\frac{\delta}{\rho-\delta},\qquad m=r+1.
\end{equation}
For example $r=\lfloor\delta/(\rho-\delta)\rfloor+1$ works. Define
\begin{equation}\label{eq:circleconstants53}
 \Delta_r=2\left(\frac{\rho r}{r+1}-\delta\right)>0,\qquad
 A_r=\frac{2(1+\rho+\delta)(2r+1)m}{\Delta_r}>1.
\end{equation}

\begin{lemma}[Explicit harmonic localization]\label{lem:harmonic53}
The nonnegative trigonometric polynomials
\[
 p_\pm(\phi)=\frac{(1\pm\cos\phi)^r}{2^{-r}\binom{2r}{r}}
\]
have Haar mean one and weighted means
\[
 \int p_\pm(\phi)\rho\cos\phi\,\frac{d\phi}{2\pi}
                         =\pm\frac{\rho r}{r+1}.
\]
For the moving feature vector
\[
 Z_m(\phi)=(\cos\phi,\sin\phi,\ldots,\cos m\phi,\sin m\phi),
\]
any terminal realization with wordwise mean error at most $\delta$ obeys
\begin{equation}\label{eq:harmonicresidual53}
 \E\norm{\E[Z_m(\phi)\mid S]}\ge
 \frac{\Delta_r}{2(1+\rho+\delta)(2r+1)}.
\end{equation}
\end{lemma}
\begin{proof}
The identity $1+\cos\phi=2\cos^2(\phi/2)$, or the constant Laurent coefficient, gives
\[
 C_r:=\int(1+\cos\phi)^r\frac{d\phi}{2\pi}
            =2^{-r}\binom{2r}{r},\qquad
 C_{r+1}/C_r=\frac{2r+1}{r+1}.
\]
Hence the weighted cosine mean is $C_{r+1}/C_r-1=r/(r+1)$. Translation by $\pi$ gives the negative case.

We use the unweighted real coefficient vector in the basis $1,\cos\phi,\sin\phi,\ldots$, without orthonormal Fourier rescaling. The absolute sum of the real cosine coefficients of either $p_\pm$, including the constant coefficient, is
\[
 B_r=\frac{4^r}{\binom{2r}{r}}\le2r+1.
\]
For $p_+$ its coefficients are nonnegative and their sum is $p_+(0)$; for $p_-$ they differ only by alternating signs. The inequality follows because the central binomial coefficient is the largest of the $2r+1$ coefficients summing to $4^r$. Multiplication by $\rho\cos\phi$ has coefficient absolute sum at most $\rho B_r$, since $\cos u\cos v=(\cos(u+v)+\cos(u-v))/2$. Thus the Euclidean norms of the moving coefficient vectors of $p_\pm$ and $p_\pm\rho\cos\phi$ are at most $2r+1$ and $\rho(2r+1)$, respectively. For completeness, let $c=\E D$, $Q=\E\norm{\E[Z_m\mid S]}$. For either weight, the conditional coefficient identity gives $|c-\int p_\pm\rho\cos\phi|\le\delta+(1+\rho+\delta)(2r+1)Q$. Subtracting the two inequalities gives \eqref{eq:harmonicresidual53}. Correctness is averaged word by word before using this identity. The starting feature norm is $\sqrt m\le m$; $A_r$ deliberately uses the latter bound to obtain a rational closed form.
\end{proof}

The integer $r$ in \eqref{eq:r53} is the localization order; it is unrelated to the number $s$ of rotation letters. For fixed subcritical error it is finite and independent of $N$ and $k$.

\begin{theorem}[A full-noise arithmetic occupation bound]\label{thm:arithmetic54}
Assume \eqref{eq:DA54} and $0\le\varepsilon<\rho/2$. Fix the localization parameters $r,m,A_r$ in \eqref{eq:r53}--\eqref{eq:circleconstants53}. For an integer $k\ge1$, put
\[
 L=\left\lceil(2k)^{1/s}\right\rceil,\qquad B=s(L-1).
\]
Every error-$\varepsilon$ realization of peak at most $k$ satisfies
\begin{equation}\label{eq:finiteDA54}
 N<B\left(1+b^{-2}(mB)^{2\tau}\log A_r\right).
\end{equation}
With no restriction on other widths, the number of cuts of width at most $k$ is at most
\begin{equation}\label{eq:DAoccupation54}
 \left\lceil B-1+B b^{-2}(mB)^{2\tau}\log A_r\right\rceil.
\end{equation}
\end{theorem}
\begin{proof}
Use $s$ consecutive groups of $L-1$ command positions. In group $i$, choose $J_i$ uniformly from $\{0,\ldots,L-1\}$, apply $J_i$ copies of the $i$th rotation, and fill the other positions with identities. This defines an external law on length-$B$ words with $L^s\ge2k$ equally likely products. On the harmonic feature use
\[
 R_m(t)=\diag(R_t,R_{2t},\ldots,R_{mt}).
\]
For two distinct count vectors $J,J'$ and any frequency $1\le l\le m$, condition~\eqref{eq:DA54} and $\|l(J-J')\|_1\le mB$ give
\[
 |e^{2\pi i l(J-J')\cdot\alpha}-1|
 =2\bigl|\sin\bigl(\pi l(J-J')\cdot\alpha\bigr)\bigr|
 \ge4b(mB)^{-\tau}=:\sigma.
\]
Here $2\sin(\pi x)\ge4x$ for $0\le x\le1/2$. Also $b\le1/2$, $m\ge2$, and $B\ge1$, so $\sigma\le1$.

For a unit vector $u=(u_1,\ldots,u_m)$ in the frequency blocks, the squared distance between the two represented orbit points is
\[
 \sum_{l=1}^m\norm{u_l}^2
       |e^{2\pi i l(J-J')\cdot\alpha}-1|^2\ge\sigma^2.
\]
An open ball of radius $\sigma/2$ about any unit center contains at most one point. At least half the $L^s$ points therefore have distance at least $\sigma/2$ from all of any $k$ centers. Since the inner-product defect between unit vectors equals half their squared distance, the finite word law has defect at least
\[
                  d_k=\sigma^2/16=b^2(mB)^{-2\tau}.
\]
Apply conditional-centroid contraction to successive blocks and Lemma~\ref{lem:harmonic53} to the terminal query. The starting norm is $\sqrt m\le m$. Thus
\[
 \left\lfloor\frac NB\right\rfloor
 \le\frac{\log A_r}{-\log(1-d_k)}
 \le b^{-2}(mB)^{2\tau}\log A_r,
\]
which proves \eqref{eq:finiteDA54}. The greedy endpoint count, with identity gaps and no prefix, gives \eqref{eq:DAoccupation54}. Internal registers remain unrestricted. Finally $B<s(2k)^{1/s}$, so \eqref{eq:finiteDA54} implies the asserted $\Omega(N^{s/(2\tau+1)})$ lower bound.
\end{proof}

\subsection{An exact polygon realization}
Write $\eta_\alpha(q)=\max_i\|q\alpha_i\|_{\R/\Z}$. The following elementary construction explains why a common denominator, rather than the number of numerical orbit points, controls the upper width.

\begin{proposition}[Exact enclosure upper bound]\label{prop:polygon54}
Let $0<\rho<1$, $a=(1+\rho)/2$, and $\gamma=\log(a/\rho)>0$. Choose an integer $q_0\ge4$ with $\cos(\pi/q_0)\ge a$. If $q\ge q_0$ and
\begin{equation}\label{eq:enclosure54}
                4\pi^2N\,\eta_\alpha(q)/q^2\le\gamma,
\end{equation}
then the indicated rotation alphabet, with or without $F$, has an exact horizon-$N$ realization with $q$ command labels.
\end{proposition}
\begin{proof}
Let $P_q$ be the regular polygon with vertices $v_j=(\cos(2\pi j/q),\sin(2\pi j/q))$, $0\le j<q$. For each $i$ choose a nearest integer $p_i$ to $q\alpha_i$ and put
\[
 t=\pi/q,\qquad h_i=2\pi|\alpha_i-p_i/q|\le t,\qquad
 \Lambda_i=\frac{\cos(t-h_i)}{\cos t}.
\]
The facet inequalities of the regular polygon give $R_{2\pi\alpha_i}P_q\subset\Lambda_iP_q$: the nearest facet normal to a rotated vertex has angular difference $t-h_i$. Further,
\[
 \log\Lambda_i
 \le\log(1+\tan(t)h_i)\le\tan(t)h_i
 \le4\pi^2\eta_\alpha(q)/q^2,
\]
since $\cos h+\tan(t)\sin h\le1+\tan(t)h$ and $\tan t\le2t$ for $0\le t\le\pi/4$. Let $\Lambda=\max(1,\Lambda_1,\ldots,\Lambda_s)$. The identity and reflection preserve $P_q$, so every command $U$ satisfies $\Lambda^{-1}UP_q\subset P_q$.

For each command and each vertex choose a probability row with barycenter $\Lambda^{-1}Uv_j$ in $P_q$. These are actual common command rows, not independent prefix/suffix factorizations. At most three successors suffice by planar convexity. Initialize the seed $\rho u_s$, $u_s\in\{\pm e_1,\pm e_2\}$, with a row whose barycenter is $au_s$. This is possible since $a\le\cos(\pi/q)$, the polygon's inradius. After a length-$N$ word the mean vertex is
\[
                         a\Lambda^{-N}U_wu_s.
\]
By \eqref{eq:enclosure54}, $\rho\Lambda^N\le a$. Hence the terminal columns
\[
                       d_j(v)=\frac{\rho\Lambda^N}{a}\,e_j^Tv
\]
belong to $[-1,1]$ and return exactly $\rho e_j^TU_wu_s$. The rows can be stationary within a horizon; the terminal scaling and chosen denominator may depend on $N$. No one finite all-horizon exact machine is asserted.
\end{proof}

\begin{lemma}[Comparable simultaneous denominators]\label{lem:denominators54}
Assume \eqref{eq:DA54}. For every integer $Q\ge2$ there is $q$ with $1\le q\le Q^s$, $\eta_\alpha(q)\le1/Q$, and
\[
 q\ge c Q^{s/(\tau+1-s)},\qquad
 c=\left(\frac{b}{s^{\tau+1}}\right)^{s/(\tau+1-s)}.
\]
In particular, for $\tau=s$ one has $cQ^s\le q\le Q^s$.
\end{lemma}
\begin{proof}
Partition the torus into $Q^s$ half-open cubes and apply the pigeonhole principle to $0,\alpha,\ldots,Q^s\alpha$. Their difference gives $q\le Q^s$ and integers $p_i$ with $|q\alpha_i-p_i|\le1/Q$.

Put $M=\lfloor q^{1/s}\rfloor$. The $(M+1)^s>q$ vectors in $\{0,\ldots,M\}^s$ cannot all have distinct residues $n\cdot p$ modulo $q$. Subtract two to obtain a nonzero vector $n$ with $q\mid n\cdot p$ and $H:=\|n\|_1\le s q^{1/s}$. Therefore
\[
 bH^{-\tau}\le\|n\cdot\alpha\|_{\R/\Z}
 \le H/(qQ),\qquad
 q^{(\tau+1-s)/s}\ge bQ/s^{\tau+1}.
\]
The displayed bound follows. This argument supplies the needed quantitative transference directly; no unspecified denominator subsequence is being used.
\end{proof}

\begin{proof}[Proof of Theorem~\ref{thm:matched54}]
The lower estimate follows from Theorem~\ref{thm:arithmetic54}. For $\tau=s$, take the constants $a,\gamma,q_0,c$ above and an integer
\[
 Q\ge\max\left(2,(q_0/c)^{1/s},
                \left(\frac{4\pi^2N}{c^2\gamma}\right)^{1/(2s+1)}\right).
\]
Choosing the least such integer, Lemma~\ref{lem:denominators54} gives $q\le Q^s=O(N^{s/(2s+1)})$, $q\ge q_0$, and
\[
 4\pi^2N\eta_\alpha(q)/q^2
 \le4\pi^2N/(c^2Q^{2s+1})\le\gamma.
\]
Proposition~\ref{prop:polygon54} supplies the exact upper bound. At and above $\rho/2$ use the fair decoder from Corollary~\ref{cor:circle53}. This proves all claims.
\end{proof}
The command count $s$ is an arithmetic parameter, not a replacement for the dimension of the physical circle. The constants deteriorate as $\varepsilon\uparrow\rho/2$ through the localization degree, and as $\rho\uparrow1$ through the enclosure slack. Both dependences are explicit in the finite inequalities. The theorem makes no uniform claim at these limiting parameters.
